\section{Розділ~\ref{sec:nora}. \nt{NORA}}

Будову \nt{NORA}-системи, її два режими, зв'язок із зіницею (не лише через \nt{NORA}), зростання відношення сигналу до фону і перевернуте U в поведінці виміряно прямо; мультиплікативне підсилення $g$ --- модель; те, що підсилення перемикає режим вибору і працює як обернена температура, --- висновки розділу; виграш від підлаштування підсилення --- розрахунок за моделлю книги; «\nt{NORA} --- ручка пошуку» і «\nt{NORA} --- переривання» --- гіпотези.

{\footnotesize
\setlength{\tabcolsep}{3pt}\renewcommand{\arraystretch}{1.15}
\begin{longtable}{>{\raggedright}p{0.5cm}>{\raggedright}p{4.0cm}>{\raggedright}p{5.6cm}>{\raggedright}p{2.3cm}>{\raggedright\arraybackslash}p{1.7cm}}
\toprule
№ & Твердження & Дослід і що виміряно & Джерело & Статус \\
\midrule\endfirsthead
\toprule
№ & Твердження & Дослід і що виміряно & Джерело & Статус \\
\midrule\endhead
1 & \nt{NORA}-нейронів у людини кілька десятків тисяч, майже всі в одній групі & Серійні зрізи стовбура мозку людини, забарвлення на тирозингідроксилазу: $53\,900$ клітин у блакитній плямі і ще $6\,260$ у сусідньому під'ядрі. & \cite{Baker1989LC} & виміряно \\
2 & Виходи розходяться майже по всьому мозку, але частини групи частково обслуговують різні області & Вірусно-генетичне мічення в мишей: \nt{NORA}-нейрони блакитної плями отримують входи з широких областей мозку і надсилають виходи в широкі області головного і спинного мозку. У щурів: група, що йде до мигдалеподібного тіла, посилює вивчення страху, а окрема група, що йде до префронтальної кори, --- його згасання. & \cite{Schwarz2015LC, Uematsu2017LC, Poe2020} & виміряно \\
3 & Фоновий режим: рівні імпульси з частотою від часток до кількох на секунду, що змінюється з рівнем неспання & Щури, запис нейронів блакитної плями в циклі сну і неспання: частота найбільша в неспанні, нижча в повільному сні, майже нуль у швидкому; зміни частоти випереджають зміну стану. & \cite{AstonJonesBloom1981} & виміряно \\
4 & Уривчастий режим: коротка пачка на важливу для задачі подію, приблизно в момент рішення & Мавпи, задача знайти рідкісну ціль: усі нейрони блакитної плями відповідають пачкою лише на ціль, через $\sim90\ms$ після неї і за $\sim200\ms$ до відповіді рукою; за поганої роботи пачки слабші. У задачі з вибором із двох пачка точніше прив'язана до відповіді, ніж до стимулу, і є як за правильних, так і за помилкових відповідей. & \cite{AstonJones1994vigilance, Clayton2004LC} & виміряно \\
5 & Діаметр зіниці --- неточна контрольна точка \nt{NORA}: він іде і за \nt{NORA}, і за \nt{ACET}, і за іншими областями & Мавпи: імпульси блакитної плями, аж до окремих імпульсів однієї клітини, передбачають розширення зіниці; схожий, але пізніший зв'язок є в горбках і поясній корі. Миші: коливання зіниці йдуть за активністю і норадреналінових, і холінергічних виходів у корі. & \cite{Joshi2016, Reimer2016} & виміряно \\
6 & Норадреналін пригнічує фонову активність сильніше за викликану; мультиплікативне підсилення~\eqref{eq:gain} --- модель, що передає цей ефект & Мавпи в стані неспання, слухова кора, іонофорез норадреналіну: фонова активність нейронів пригнічується сильніше, ніж відповідь на голоси родичів, --- відношення сигналу до шуму зростає. Мультиплікативну сигмоїду~\eqref{eq:gain} взято з мережевої моделі; буквального множення крутизни не виміряно. & \cite{Foote1975NE, ServanSchreiber1990} & виміряно; $g$ --- модель \\
7 & Підсилення підвищує $d'$ лише проти шуму після підсилювача~\eqref{eq:gainsnr} (твердження~\ref{prop:gainnoise}) & Виведено в розділі; у мережевій моделі підсилення покращує виявлення сигналу. Досліду, який відокремлював би шум до і після підсилювача і перевіряв би цю різницю, немає. & виведено в розділі; \cite{ServanSchreiber1990} & теорія \\
8 & Межа $g\beta=1$ перемикає мережу вибору з «розмірковує» на «вирішила»~\eqref{eq:wtagain} (твердження~\ref{prop:gainwta}, рис.~\ref{fig:pitchfork}) & Лінійний аналіз двох груп, що гальмують одна одну; виведено в розділі. Прямих вимірювань цього порога в мозку немає. & виведено в розділі & теорія \\
9 & Коефіцієнт підсилення --- обернена температура вибору~\eqref{eq:softmax} & За логістичного шуму після підсилювача ймовірність вибору --- softmax з $T=\sigma/g$; виведено в розділі. & виведено в розділі & теорія \\
10 & \nt{NORA} задає, скільки шукати: висока фонова активність --- мале діюче $g$ (пошук), уривчаста пачка в момент рішення --- велике (рішення) & Люди: перед вибором навмання базова зіниця ширша, ніж перед вибором найкращого відомого; у людей із ширшою зіницею схильність шукати вища. Щури: перехід у режим випадкового вибору керується входом \nt{NORA} у передню поясну кору. Що пачка підвищує діюче $g$, прямо не виміряно. & \cite{JepmaNieuwenhuis2011, Tervo2014stochastic, AstonJones2005} & гіпотеза \\
11 & Винагорода залежить від $g$ як перевернуте U; найкраще $g$ менше у світі, що швидко змінюється (твердження~\ref{prop:explore}, рис.~\ref{fig:invertedU}) & \texttt{models/nora\_explore.py}, $60$ прогонів по $4000$ спроб. Зміна раз на $1000$ спроб: найкраще $g=16$, частка $0.976$; раз на $20$: $g=8$, частка $0.886$; при $g=0.5$ --- $0.77$. Перерахунок збігся з розділом. & \texttt{models/\allowbreak nora\_\allowbreak explore.py} & модель \\
12 & Перевернуте U в поведінці: найкраще за помірної фонової активності з чіткими пачками & Мавпи, та сама задача, що в рядку~4: у періоди доброї роботи фонова частота помірна, а пачки на ціль чіткі; за високої фонової частоти пачок немає і хибних відповідей більше. Зміни фонової і викликаної активності йдуть за коливаннями якості роботи. & \cite{Usher1999LC, AstonJones2005} & виміряно \\
13 & \nt{NORA} повідомляє неочікувану невизначеність, \nt{ACET} --- очікувану & Теорія байєсівського висновку з двома видами невизначеності; прямого досліду, що розділяє ці сигнали за медіаторами, у цитованих роботах немає. & \cite{YuDayan2005} & гіпотеза \\
14 & Після різкої зміни люди навчаються швидше, і зіниця розширюється & Люди, задача передбачення з раптовими змінами: короткі зміни зіниці йдуть за оцінкою ймовірності зміни і разом із базовою зіницею передбачають, наскільки сильно нові дані змінюють оцінку (швидкість навчання); зміна зіниці, не пов'язана зі світлом і задачею, змінює цю швидкість. Що зіниця тут показує саме \nt{NORA}, --- висновок за непрямими даними рядка~5. & \cite{Nassar2012} & виміряно \\
15 & Уривчаста пачка працює як переривання: мережа скидає стан & Прямого досліду, де пачка \nt{NORA} скидає стан мережі, немає; виміряно лише пачки на важливі події (рядок~4). & \cite{BouretSara2005} & гіпотеза \\
16 & Підлаштування $g$ або швидкості навчання за несподіваністю майже не краще за найкращі сталі налаштування & \texttt{models/nora\_explore.py}, світ з епохами по $1000$ спроб (зміна раз на $1000$ і раз на $20$). Найкраще стале $g=8$: $0.925$; підлаштування $g$: до $0.927$; підлаштування швидкості навчання: до $0.926$; стала швидкість $0.4$: $0.928$. Перерахунок збігся з розділом. & \texttt{models/\allowbreak nora\_\allowbreak explore.py} & модель \\
\bottomrule
\end{longtable}}
